\section{Removing the entire cubic coefficient-two image}
\label{sec:fr-cubic-coefficient-two}

Use \(M=\ker\lambda\), \(D=N_{\mathrm{rad}}\), and
\[
 R_D=D\cap\ker\Phi_2,\qquad G=R_DH.
\]
The corrected cubic frame \(T\) is the homomorphism of
Lemma~\ref{lem:adelic-corrected-frames}. We distinguish its
\emph{coefficient-two component} from the maximal pro-two factor
of its whole adelic target. Put
\[
 O=\mathbf Z_2[\omega],\qquad
 \delta=1+2\omega,\qquad \delta^2=-3,
\]
\[
 \mathcal U=\{z\in1+4O:z\bar z=1\},\qquad
 \mathfrak t_2:M\longrightarrow\mathcal U,\quad
 \mathfrak t_2(g)=T_{g,2}.
\]
The corrected frame already belongs to \(1+4O\); no additional
projection away from roots of unity is needed in this definition.

\begin{lemma}\label{lem:fr-cubic-two-character}
The logarithm identifies \(\mathcal U\) with
\(4\delta\mathbf Z_2\). The formula
\[
 \tau(g)=\frac{\log\mathfrak t_2(g)}{4\delta}\pmod2
 \qquad(g\in M)
\]
defines a continuous homomorphism \(M\to\mathbf F_2\), invariant
under conjugation by all of \(G\). Equivalently,
\[
 \mathfrak t_2(g)\equiv(1+4\tau(g))I\pmod8.
\]
Its kernel in \(\mathcal U\) is \(\mathcal U^2\). An element of
\(\mathcal U\) with \(\tau=1\) topologically generates \(\mathcal U\).
\end{lemma}

\begin{proof}
Logarithm and exponential are inverse continuous homomorphisms
between \(1+4O\) and \(4O\). The norm condition becomes trace zero:
\(\log(z\bar z)=\operatorname{Tr}\log z\).
Since
\[
 \operatorname{Tr}(a+b\omega)=2a-b,
\]
the trace-zero lattice in \(O\) is precisely
\(\delta\mathbf Z_2\). Hence
\(\log\mathcal U=4\delta\mathbf Z_2\).
This proves the assertions concerning the logarithm, squares, and
topological generators.

Modulo eight, exponential is \(x\mapsto1+x\), and
\(\delta\equiv1\pmod2\). This gives the displayed scalar formula.
The frame is a homomorphism on \(M\), so \(\tau\) is a continuous
homomorphism. The integral conjugation formula
\eqref{eq:adelic-cubic-kernel-conjugation} identifies
\(T_{ghg^{-1},2}\), for \(h\in M\), with an integral similarity
of \(T_{h,2}\). Its reduction modulo eight fixes a scalar matrix.
Thus \(\tau(ghg^{-1})=\tau(h)\) for every \(g\in G\).
The same formula makes the kernel of \(\mathfrak t_2\) in \(M\)
normal in \(G\).
\end{proof}

\begin{lemma}\label{lem:fr-cubic-two-arithmetic-bit}
The restriction of \(\tau\) to \(D\cap H\) is nonzero. There is,
on the other hand, an element \(k_0\in D\cap H\) with
\[
                  \eta_3(k_0)=1,\qquad \tau(k_0)=0.
\]
\end{lemma}

\begin{proof}
Let
\[
 k=\mathbf Q^{\mathrm{ab}},\quad
 \alpha=3^{1/4},\quad
 v=(2+\sqrt3)^{1/4},\quad
 L_{\mathrm{rad}}=k(q^{1/n}:q\text{ prime},\ n\ge1).
\]
The established eight-division calculation
(Lemma~\ref{lem:lr-cubic-eight-field}) gives
\[
 k(E_3[8])=k(\alpha,v),\qquad [k(\alpha,v):k]=4,
\]
and identifies the action on \(v\) with the corrected scalar bit
\(\tau\). In particular \(v\in K_3^{\mathrm{ab}}\), and
\(v\notin k(\alpha)\). Also \(k(v)/k\) is quadratic, since
\[
 v^2=\frac{\sqrt6+\sqrt2}{2}\in k.
\]

We claim that \(v\notin L_{\mathrm{rad}}\).
The exact radical--CM intersection theorem gives
\[
 I=L_{\mathrm{rad}}\cap
       (K_3^{\mathrm{ab}}K_4^{\mathrm{ab}})
   =k(q^{1/12}:q\text{ prime},\,2^{1/8},3^{1/8}).
\]
Its explicit product Galois group shows that its largest elementary
abelian subextension of exponent two is
\[
                         k(q^{1/4}:q\text{ prime}).
\]
This follows coordinate by coordinate from the cyclic factors of
orders twelve and six; every continuous quadratic character uses
only finitely many coordinates. Each rational fourth root is abelian
over \(K_4=\mathbf Q(i)\), so this elementary subextension is contained
in \(K_4^{\mathrm{ab}}\).
If \(v\) belonged to \(L_{\mathrm{rad}}\), its quadratic field over
\(k\) would consequently lie in both CM fields. But
Lemma~\ref{lem:adelic-cm-entanglement} gives
\[
 K_3^{\mathrm{ab}}\cap K_4^{\mathrm{ab}}=k(\alpha),
\]
contrary to \(v\notin k(\alpha)\).
Thus \(L_{\mathrm{rad}}(v)/L_{\mathrm{rad}}\) is quadratic.
An element of \(G_{L_{\mathrm{rad}}}=D\cap H\) moving \(v\)
has \(\tau=1\).

For the second assertion, choose the established arithmetic element
\(\sigma\in M\cap H\) with
\[
 \rho_2^I(\sigma)=2,\qquad
 \rho_q^I(\sigma)=0\quad(q\text{ odd}).
\]
The cubic congruence
\(\rho_2^I\equiv-2m_3\pmod4\) makes \(m_3(\sigma)\) odd.
Choose \(h\in H\) with \(\lambda_2(h)=3\), and write
\(t=\lambda(h)\in\widehat{\mathbf Z}^{\times}\). Set
\[
                 k_0=(h\sigma h^{-1})\sigma^{-t}.
\]
The exponent here is the full profinite exponent \(t\), not the
integer three. Exact weight-one conjugation of every rational
Kummer cocycle gives \(\rho_q^I(k_0)=0\), hence \(k_0\in D\cap H\).
The cubic scalar law gives
\[
 m_3(k_0)=(3^3-3)m_3(\sigma)=24m_3(\sigma),
\]
so \(\eta_3(k_0)=1\).
By Lemma~\ref{lem:fr-cubic-two-character}, \(\tau\) is a
\(G\)-invariant homomorphism on all of \(M\). Since \(t_2\) is odd,
\[
 \tau(k_0)=\tau(\sigma)-(t_2\bmod2)\tau(\sigma)=0.
\]
\end{proof}

\begin{theorem}\label{thm:fr-cubic-two-exact-image}
The full and arithmetic coefficient-two images are
\[
 \boxed{\quad
 \mathfrak t_2(R_D)
   =\mathfrak t_2(R_D\cap H)
   =\mathcal U.
 \quad}
\]
In particular the upper image is the entire level-four norm-one
group, rather than just its square subgroup.
For
\[
                    R'=R_D\cap\ker\mathfrak t_2
\]
one has
\[
 R_D/R'\simeq\mathcal U\simeq\mathbf Z_2,\qquad
 G=R'H,\qquad
 B\simeq R'/\overline{[G,R'](R'\cap H)}.
\]
Thus \(R'\) is a proper closed normal subgroup of infinite index in
\(R_D\), and it preserves the full relative abelian obstruction.
\end{theorem}

\begin{proof}
Write \(J=D\cap H\), \(J_\tau=J\cap\ker\tau\), and
\(P_0=\Phi_2(D)\). The established dyadic normalization gives
\[
 \Phi_2(J)=P_0,\qquad
 H^1(P_0,\mathbf F_2)^G=\mathbf F_2\eta_3
\]
by Corollary~\ref{cor:protwo-image-invariant-characters}.
Let \(L=\Phi_2(J_\tau)\).
It is compact, hence closed, and normal in \(P_0\).
The quotient \(P_0/L\) is a quotient of \(J/J_\tau\); its order
is at most two.

Moreover \(L\) is stable under \(H\), since \(J_\tau\) is.
The action of \(G\) on \(P_0\) factors through its dyadic realization.
The factorization \(G=R_DH\), with \(\Phi_2(R_D)=1\), therefore
makes \(L\) stable under all \(G\).
If \(L\ne P_0\), the nonzero character of its quotient is
\(G\)-invariant and so equals \(\eta_3\).
But the element \(k_0\) of the preceding lemma belongs to
\(J_\tau\) and has \(\eta_3(k_0)=1\), a contradiction.
Consequently
\[
                         \Phi_2(J_\tau)=P_0.
\]

Choose \(r\in J\) with \(\tau(r)=1\).
There is \(y\in J_\tau\) with
\(\Phi_2(y)=\Phi_2(r)\). Then
\[
                         w=ry^{-1}\in R_D\cap H,
 \qquad \tau(w)=1.
\]
Its coefficient-two value topologically generates \(\mathcal U\)
by Lemma~\ref{lem:fr-cubic-two-character}.
The compact procyclic closure of \(w\) therefore maps onto
\(\mathcal U\). This proves arithmetic surjectivity; the full image
is contained in \(\mathcal U\) by definition, proving equality.

The kernel \(R'\) is closed and \(G\)-normal. For every \(r\in R_D\),
choose \(j\in R_D\cap H\) with
\(\mathfrak t_2(j)=\mathfrak t_2(r)\). Then \(rj^{-1}\in R'\).
Thus
\[
 R_D=R'(R_D\cap H),\qquad G=R'H.
\]
The compact quotient theorem gives \(R_D/R'\simeq\mathcal U\).

Finally, for any finite abelian group \(A\), a \(G\)-invariant
character of \(R'\) killing \(R'\cap H\) extends uniquely to
\(G\) by \(\chi(rh)=\chi(r)\). This is well-defined and
continuous by the compact quotient map \(R'\times H\to G\).
It is exactly a character of \(G\) trivial on \(H\), equivalently
a character of the established relative abelian obstruction \(B\).
Character separation over all finite abelian targets gives the
displayed canonical presentation of \(B\).
\end{proof}

\begin{corollary}\label{cor:fr-cubic-two-joint-gate}
Let
\[
 V_D=
 \bigl(D\cap\bigcap_p\ker\Phi_p\cap\ker d_{(3)}\bigr)
       \cap\ker W_{\mathrm{odd}}\cap\ker T_{\mathrm{odd}},
 \qquad
 V'=V_D\cap\ker\mathfrak t_2.
\]
Then
\[
 \mathfrak t_2(V_D)=\mathfrak t_2(V_D\cap H)=\mathcal U,
 \qquad [V_D:V']=[V_D\cap H:V'\cap H]=\infty,
\]
and
\[
 B_{\max2}\simeq
 \left(V'/\overline{[G,V'](V'\cap H)}\right)_{\max2}.
\]
If \(X_2=H^1(N,\mathbf F_2)^G\), its exact restriction kernels
to \(R'\) and to \(V'\) are both
\[
 \boxed{\quad
 \left(\bigoplus_{q\ \mathrm{odd}}\mathbf F_2\psi_q\right)
 \oplus\mathbf F_2\eta_3\oplus\mathbf F_2(\tau|_N).
 \quad}
\]
The last summand is independent of all preceding summands.
\end{corollary}

\begin{proof}
The established reductions show that \(R_D/V_D\) is pro-odd.
Explicitly, \(R_D/P_D\) is pro-odd; the successive cubic-defect,
odd harmonic, and odd cubic quotients are pro-odd as well.
Extensions of pro-odd groups are pro-odd.

For the element \(w\in R_D\cap H\) constructed above, take
\(w_2=w^{\varepsilon_2}\), where
\(\varepsilon_2\in\widehat{\mathbf Z}\) is one at two and
zero at all odd primes. Its image in the pro-odd quotient
\(R_D/V_D\) is one, so \(w_2\in V_D\cap H\).
Its image in the pro-two group \(\mathcal U\) is unchanged.
It therefore still topologically generates \(\mathcal U\), proving
the exact images and both infinite indices.

The quotient
\[
                         R'/V'\hookrightarrow R_D/V_D
\]
is pro-odd. Apply the relative pro-odd restriction lemma to the
already established factorization \(G=R'H\).
This gives the asserted presentation of \(B_{\max2}\).
The factorization \(G=R'H\) is essential to this application;
an internal arithmetic factorization inside \(V_D\) would not
by itself give the conclusion.

For the character kernel, suppose \(\chi\in X_2\) vanishes on
\(R'\). Its restriction to \(R_D\) factors through
\(R_D/R'\simeq\mathcal U\). The continuous quadratic character
space of \(\mathcal U\simeq\mathbf Z_2\) is generated by \(\tau\).
Subtracting either zero or \(\tau|_N\) therefore gives a character
vanishing on \(R_D\). The known kernel there is precisely the
Kummer span plus \(\eta_3\).
Conversely all displayed characters vanish on \(R'\).
The new summand is independent because \(\tau(w)=1\), whereas
all the older kernel characters vanish on \(R_D\).

If \(\chi\) vanishes on \(V'\), its restriction to \(R'\) factors
through the pro-odd quotient \(R'/V'\), hence is zero.
Thus restriction to \(V'\) has the same kernel.
\end{proof}

\subsection{The arithmetic class of the new quadratic lift}

We identify \(\tau|_{G_{K_2}}\) in the explicit arithmetic target
of Theorem~\ref{thm:fr-Y2-complete}.
Let \(a,b,c_3\) denote the quadratic cyclotomic characters of
\(\sqrt2,\sqrt{-1},\sqrt{-3}\), respectively, and let
\(\delta_2=b^2\), with \(\delta_3\) the cyclic lift obstruction
at the odd prime three. Products below are cup products.
The target normalization is
\[
 \operatorname{tr}(\kappa_{2,3})
  =\operatorname{Inf}(a c_3+\delta_2+\delta_3),\qquad
 \operatorname{tr}(\psi_3)
  =\operatorname{Inf}(b c_3+\delta_3).
\]

\begin{proposition}\label{prop:fr-cubic-two-Y2-class}
One has the exact equality of arithmetic characters
\[
             \tau|_{G_{K_2}}=\kappa_{2,3}+\psi_3.
\]
Consequently
\[
                \vartheta_{2,3}:=\tau|_N+\psi_3\in X_2
\]
is an explicit invariant lift of the corrected cross-cup target
\(\kappa_{2,3}\).
The last summand in Corollary~\ref{cor:fr-cubic-two-joint-gate}
can equivalently be written \(\mathbf F_2\vartheta_{2,3}\).
Neither character is asserted to be arithmetic-trivial.
\end{proposition}

\begin{proof}
Put \(L_v=\mathbf Q(i,v)\), where \(v=(2+\sqrt3)^{1/4}\).
The number
\[
 u=v^2=(\sqrt6+\sqrt2)/2
\]
has degree four over \(\mathbf Q\), since
\(u-u^{-1}=\sqrt2\) and \(u^2-2=\sqrt3\).
As \(v\notin k\), adjoining its square root to
\(\mathbf Q(u)\subset k\) has degree two. Thus
\([\mathbf Q(v):\mathbf Q]=8\), and this real field does not
contain \(i\), giving \([L_v:\mathbf Q]=16\).
All conjugates of \(v\) are \(i^jv^{\pm1}\), so \(L_v/\mathbf Q\)
is Galois.

Its automorphisms have the unique forms
\[
 g=(j,e,b):\qquad
 v\longmapsto i^jv^{(-1)^e},\quad
 i\longmapsto i^{(-1)^b},
 \qquad j\in\mathbf Z/4,\quad e,b\in\mathbf F_2.
\]
Their multiplication is
\[
 (j,e,b)(j',e',b')
   =\bigl((-1)^b j'+(-1)^{e'}j,\ e+e',\ b+b'\bigr).
\]
The element \(z=(2,0,0)\) is central and negates \(v\).
Its fixed field is
\[
                   L_0=\mathbf Q(i,\sqrt2,\sqrt3),
\]
of degree eight. The characters of this quotient on the three
square roots \(\sqrt2,\sqrt{-1},\sqrt{-3}\) have values
\[
                         A=j+e,\qquad b,\qquad C=b+e
                         \quad\text{in }\mathbf F_2.
\]
In particular \(e=b+C\), and \(j\bmod2=\alpha=A+b+C\).

Choose the section of this elementary abelian quotient with
\(j=\alpha\in\{0,1\}\).
The central carry, obtained from the multiplication law, is
\[
 \frac{(-1)^b\alpha'+(-1)^{e'}\alpha
                    -(\alpha+\alpha'\bmod2)}2
  =\alpha\alpha'+b\alpha'+\alpha e'\pmod2.
\]
Consequently the class of the central extension is
\[
                         \alpha^2+b\alpha+\alpha e
                  =a^2+b^2+a c_3+b c_3
\]
in the cohomology of the three-character quotient.
For an explicit check of the cup-order change, expansion gives
the cocycle \(AA'+CA'+bb'+bC'\); replacing \(CA'\) by \(AC'\)
changes it by the coboundary of the function \(AC\).

Pull this extension back to the full cyclotomic unit group
\(U=\widehat{\mathbf Z}^{\times}\).
The character \(a\) of \(\sqrt2\) lies on the procyclic factor
\(1+4\mathbf Z_2\) and lifts to a \(\mathbf Z/4\)-valued
character. Hence \(a^2=0\) in \(H^2(U,\mathbf F_2)\).
The pulled-back class is therefore
\[
                            b^2+a c_3+b c_3.
\]
The quadratic character of \(G_k\) defined by this central
extension is its action on \(v\), hence is \(\tau\) by the
eight-division comparison. Passing to \(G_{K_2}\) inflates the
same class from \(U\) to
\(\operatorname{Gal}(K_2/\mathbf Q)\).
The transgression used in the arithmetic-target theorem is
injective. Finally,
\[
 (a c_3+\delta_2+\delta_3)+(b c_3+\delta_3)
                          =a c_3+b c_3+b^2.
\]
This proves the claimed character identity, with no uncomputed
diagonal or cubic-edge correction.
The stated lift and the equivalent kernel expression follow.
\end{proof}

\paragraph{Scope of the normalization.}
The eliminated group is the coefficient-two level-four norm-one
torus \(\mathcal U\). The finite two-Sylow subgroups at odd coefficient
primes are not included in this assertion. The proof establishes
the full arithmetic matching on \(R_D\) before passing to \(V'\).
It neither identifies the remaining relative quotient with zero
nor asserts \(G=V'H\).
